\section{Behavior Under Algorithmic Redistricting}
\label{sec:behavior}

\subsection{Why Bisect Achieves High Responsiveness}

Responsiveness $R \approx 2.0$ under \textsc{bisect} reflects the distribution
of district vote shares that compactness-driven redistricting produces. The
$S(v)$ curve is steep near 50\% when many districts have vote shares close to
50\%---competitive districts that will flip at small swings. It is flat near
50\% when few districts are close to 50\%, because the swing needed to flip any
one of the distant safe seats is large.

Under \textsc{bisect}, compact districts centered on geographically mixed areas
tend to have vote shares that cluster near the statewide median. This is not a
designed property---the algorithm does not target any particular vote-share
distribution---but an emergent consequence of sampling populations from
geographically contiguous areas in states where partisan communities are
geographically mixed at fine scales. The result is a roughly bell-shaped
distribution of district vote shares centered near the statewide mean, with
the tails (safe seats) under-represented relative to a plan that deliberately
concentrates partisan voters.

When the district vote-share distribution is roughly uniform or bell-shaped near
50\%, many districts cluster near the pivot, and $R$ is high. When the distribution
is bimodal (many safe seats far from 50\%, few competitive districts near 50\%),
$R$ is low because the curve is flat near 50\% and steep only in the vote-share
ranges where safe seats cluster.

\subsection{Why Enacted Maps Have Suppressed Responsiveness}

Enacted maps with deliberate partisan gerrymandering produce bimodal district
vote-share distributions: a cluster of safe Republican seats (vote shares near
35--42\%) from cracking, and a cluster of packed Democratic seats (vote shares
near 70--80\%), with few competitive districts near 50\%. This bimodal
distribution produces a flat $S(v)$ near 50\%: the curve jumps at the cracked-
district vote shares (around 35--42\% Democratic, which flip at swings of
$0.50 - 0.42 = 0.08$ or larger) and at the packed-district vote shares
(around 70--80\%, which flip only at very large Democratic waves), with a
flat region in between where no districts are available to flip.

The flat $S(v)$ near 50\% translates to $R < 1.5$: each point of statewide
vote share translates to fewer than 1.5 seat-share points near the pivot.
For NC ($k=14$) with enacted $R \approx 1.3$, a 3.5-point Democratic swing
moves the expected seat share by approximately $3.5 \times 1.3 / 100 \approx
0.046$, i.e., 0.65 seats. Republican incumbents in cracked-but-safe seats
are insulated from all but extraordinary wave elections.

\subsection{The Graphical Exhibit}

The most informative way to present the seats-votes curve as courtroom evidence
is as a comparison graph:
\begin{itemize}
\item Horizontal axis: statewide Democratic two-party vote share $v \in [0.30, 0.70]$
  (the plausible range for NC).
\item Vertical axis: expected Democratic seat share $S(v) \in [0, 1]$.
\item Plot \textsc{bisect} $S(v)$: a step function passing through $(0.50, 0.50)$
  with slope $\approx 2.0$ and approximately symmetric shape.
\item Plot enacted $S(v)$: a step function passing below $(0.50, 0.50)$---at
  $v = 0.50$, enacted $S(0.50) \approx 0.43 = 6/14$---with slope $\approx 1.3$
  and an asymmetric shape that favors Republicans across the full range.
\item Mark the reference point $(0.50, 0.50)$ and the anti-symmetric ideal
  for reference.
\end{itemize}
This graph communicates two properties simultaneously without any scalar calculation:
(1) the vertical gap between the bisect and enacted curves at $v=0.50$ is the
Partisan Bias difference; (2) the relative steepness of the two curves at $v=0.50$
is the responsiveness difference. Courts can read both properties directly from
the graph.
